\chapter{结论与讨论}\label{chap4:discussion}

% ─────────────────────────────────────────────
\section{结论}\label{chap4_1:conclusion}
% ─────────────────────────────────────────────

本研究的核心贡献是在猕猴研究中将财富概念引入传统的经济决策框架，为理解财富的神经表征及其对决策过程的调控机制提供了新的视角和证据。

赚取–消费范式在两方面弥补了猕猴研究中传统代币任务的不足：其一，取消了固定的代币兑换阈值，避免了外生参考点对财富效应解释的混淆；其二，将选择置于赚取（次级强化物）与消费（初级强化物）之间，构成了跨类别的价值权衡——这是日常经济决策的基本形式之一，却鲜少在动物电生理研究中被直接考察。据我们所知，这是首次在非人灵长类动物中系统考察财富依赖的赚取–消费决策并同步记录前额叶皮层电生理活动的研究。

猕猴在赚取–消费范式中的决策偏好与其财富水平密切相关：随着财富的累积，消费倾向逐渐上升。这一模式与边际效用递减原理的定量预测相符，并可由效用模型以四个自由参数加以描述。模型的关键参数$\alpha$显著小于1，直接验证了凹效用函数的存在。这些结果确立了赚取–消费范式在灵长类中的有效性，也为后续神经分析提供了可计算的选项价值度量。

在神经水平上，财富依赖的经济决策相关机制有三个核心发现。第一，三个脑区中均存在对当前财富状态的稳定表征，其中ACC在选项呈现阶段对财富的编码尤为突出。第二，ACC和OFC精确追踪了财富状态的变化；DLPFC虽也编码财富水平，但并未参与这一动态更新。第三，上述财富表征系统性地调制了前额叶对选项价值的编码，进而影响后续的经济决策。

% ─────────────────────────────────────────────
\section{讨论}\label{chap4_2:discussion}
% ─────────────────────────────────────────────

\subsection{货币的通用媒介属性}\label{chap4_2_1:serial_parallel}
货币的核心属性之一是充当一般等价物：任何商品或劳动，无论其具体形态如何，均可通过货币折算为统一的价值量纲，从而实现跨类别的比较与交换。本实验中的"财富"变量在功能上扮演了类似角色，将付出与获取统一表达为可比的神经信号。然而，当前设计在消费和赚取两个维度上均较为单一，这削弱了对本研究中的“财富”的通用媒介属性。未来可将消费选项扩展至不同种类的果汁\citep{padoaschioppa2006neurons}、社会性图片\citep{deaner2005monkeys}等，将赚取选项扩展至按压杠杆\citep{kennerley2009neurons}等不同类型的动作，从而更全面地评估财富的神经表征是否具有跨类别的稳定性。

\subsection{物品域–动作域决策}\label{chap4_2_2:serial_parallel}

Padoa-Schioppa\citep{padoaschioppa2011neurobiology}提出了串行计算的经济决策模型：物品域（good-based）主观价值的计算与比较在OFC完成，决策结果随后由其他脑区（如ACC，lPFC）转换为动作域信号。然而，一些研究对其提出了挑战。Lin等人\citep{lin2020evidence}发现DLPFC同时编码物品域和动作域信息，且物品域信号并不早于动作域信号；Hunt等人\citep{hunt2015capturing}则发现DLPFC、OFC和ACC三个脑区内的价值相关动态几乎同步展开。这些结果表明，前额叶不同亚区可能在各自的参考框架内并行处理决策信息，而非严格按照串行模型中的先后顺序进行转换。

图~\ref{fig3_5:latency}的结果显示OFC的物品域信号在时间上并不早于DLPFC的动作域信号，与串行计算的假设不符；另一方面，虽然OFC和DLPFC分别偏好物品域和动作域信号，但其非偏好域的信号依然可以被检测到，且潜伏期晚于偏好域信号，单纯的并行模型也无法解释这一现象，提示脑区间有更为复杂的交互模式。Rushworth等人\citep{rushworth2012valuation}提出的折中视角或许更为贴切：前额叶并不存在单一通用的决策回路，而是包含多个串行与并行交织的系统，不应以简单的二分法来理解其决策机制。

\subsection{选择切换}\label{chap4_2_3_switch:switch}

在当前任务中，从赚取转向消费（或反之）的每一次切换，在物品域上意味着选项类别的改变（代币vs.饮用水），在动作域上意味着眼跳方向的改变——两个维度完全共线，无法独立分析切换决策的神经基础。

Hayden等人\citep{hayden2011neuronal}提供了一种解耦思路：每个试次均需执行具体的动作，选择历史自然决定了"留守（stay）"与"切换（switch）"而不需要额外报告。将类似逻辑引入赚取–消费范式，或可更清晰地研究转换决策本身的神经机制。

\subsection{财富的数值型编码与数量感知}\label{chap4_2_4:numerosity}

我们在实验中发现了以数值形式编码财富水平的神经元（图~\ref{fig4_1:example}D），其调谐曲线以偏好财富水平为峰值、向两侧递减，这与Yang等人\citep{yang2022primate}的结果类似。在他们的结果中，数值型神经元所占比例（38\%）远高于参数型（5\%）和类别型（5\%）神经元。说明数值型神经元是财富编码的主要形式。这些发现均可在数量感知（numerosity）的框架中得到解释。Nieder等人\citep{nieder2002representation}首次发现猕猴外侧前额叶皮层（lateral prefrontal cortex，lPFC）神经元对物体数量的单峰调谐。Nieder等人\citep{nieder2003coding}进一步发现，这些调谐曲线在对数尺度而非线性尺度上呈现近似对称的形态，说明数量编码符合Weber–Fechner定律。Nieder等人\citep{nieder2007code}将上述编码方案扩展至1–30的数字范围也得到了一致的结论。同时，在他们的结果中，偏好各数字的神经元比例呈U形分布，与本研究偏好财富水平分布的U形形态（图~\ref{fig4_2:proportion}）十分相似。代币数量作为一种离散的数值量，其神经机制可能与视觉数量感知相同。

但在当前任务中，所有代币的面值相同，导致财富水平与代币数量无法区分，无法验证财富编码是否可以作为抽象的数量独立于视觉信号。引入不同面值的代币（如不同颜色、不同形状的图形）可以在实验设计上解耦财富水平与代币数量，使得同一财富水平下的不同代币组合成为可能。这可以帮助我们排除代币数量这一替代解释，并检验财富作为抽象数量，其神经表征是否与真实的视觉数量采用相同的机制。

同时，本研究观察到的U形分布可能是实验设计与分析方法造成的伪迹。以钱包上限（25）为例，任何偏好财富水平大于等于25的神经元在当前分析框架中都会被截断并归类为偏好25，导致高端区间的神经元数量被高估。即我们无法区分单调型神经元和峰值在钱包上下限之外的数值型神经元。扩展钱包容量的上下限可能可以检验这一假设。若分布保持稳定，则代表了神经元对特定财富水平的真实偏好；反之，则支持截断效应的假设。需要注意的是，猕猴在实验中可能会适应钱包上下限的变化，U形分布有可能会跟随之改变。如果上述讨论为真，这种适应在神经元层面如何发生是值得进一步研究的问题。

\subsection{真实参考点}\label{chap4_2_5:reference_point}

即使本研究在任务设计上去除了固定兑换阈值这一外设参考点，猕猴实际使用的内部参考点仍是一个无法直接观测的未知量。财富神经元编码的究竟是代币的绝对数量，还是相对于某个内部参考点的偏差，现有数据无法给出明确答案。Ferro等人\citep{ferro2026accumulation}通过预设参考点与实际财富水平的依赖关系，提供了一种从行为数据估计参考点位置的建模方法。但其适用范围还有待进一步验证。未来研究可从两方面入手：首先是将上述建模方法扩展到当前任务，其次是尝试在实验设计中让被试主动报告所使用的参考点。

\subsection{财富表征的操控}\label{chap4_2_6:causal}

当前研究提供的全部证据均为相关性：我们知道这些脑区编码财富，但尚不能证明这些编码信号是动物产生财富依赖行为的必要条件。一个核心问题是：能否通过操控这些脑区的神经活动，改变动物对当前财富状态的主观估计，进而系统地影响其决策？

这一问题面临一个根本挑战：财富与价值在脑区和神经元水平均无法完全分离。同一脑区中编码财富的神经元可能也参与价值计算，因此脑区粒度的操纵将不可避免地同时干扰财富表征与价值编码，使两者的因果贡献无法独立评估。

Ballesta等人\citep{ballesta2020values,ballesta2022orbitofrontal}在OFC的工作提供了新的方法论视角：将不同信息的呈现在时间上精确分离，并在特定时间窗口内施加低电流刺激，可以在尽量不触及另一信息编码的前提下选择性干扰目标信号。将这一思路移植至财富–价值分离的场景：若能在任务设计中先呈现财富状态，间隔一段时间后再呈现该试次的选项（或反之），则可在财富编码活跃而价值编码尚未启动的时间窗口内实施扰动，从而实现对财富表征的相对选择性操控，并观察其对动物决策行为的影响。

\subsection{财富水平对价值编码的影响}\label{chap4_2_7:timing}

图~\ref{fig5_3:model}显示，财富水平对价值编码的调制效应仅在钱包更新阶段存在，而在选项刚呈现时比较微弱。

这一结果可能受限于分析的灵敏程度。在当前任务中，$V_E$与固定代币水平下的赚取价值$V_E^{\mathrm{fixed}}$之间的相关性高达约0.8，这限制了回归模型对两个变量的区分度。若能在实验设计中扩展财富的变化范围，使不同财富水平下的$V_E$与$V_E^{\mathrm{fixed}}$有更大差异，则有望进一步检验财富水平对价值编码的调制作用。

另一方面，若三个脑区在选项呈现期的财富调制确实不显著，则与行为结果产生矛盾：行为上财富对选择的影响十分明显。这将提示财富–价值整合可能并不发生在本研究记录的三个区域，而是在其他脑区中完成后，再将整合结果以价值信号的形式传至前额叶。

\subsection{价值比较的候选脑区}\label{chap4_2_8:deltaV}

本研究在三个脑区均未发现显著的价值差（$\Delta V$）编码，提示显式的价值比较信号可能不在PFC内部完成。

腹侧纹状体（VS）是一个可能的候选脑区。Samejima等人\citep{samejima2005representation}在猕猴纹状体发现了部分神经元编码了两个动作的价值之差。Strait等人\citep{strait2015signatures}比较了VS与vmPFC的编码模式，发现VS神经元同样编码选项价值、价值差以及被选价值，且VS中被选选项的表征时间早于vmPFC中的对应信号——这提示VS可能主动参与价值比较过程，而非仅接收皮层的价值计算输出。

\subsection{OFC}\label{chap4_2_9:OFC}

本研究中OFC的发现在经济决策方面与Padoa-Schioppa等人建立的经典体系高度一致。在该框架下，OFC神经元以独立于空间位置和运动规划的方式编码决策所需的内部变量，分别对应选项提供价值（offer value）、被选选项价值（chosen value）和选择结果（choice）\citep{padoaschioppa2006neurons,padoaschioppa2017orbitofrontal}。本研究观察到的$V_E$、$V_S$和$V_\mathrm{Chosen}$编码结构与这一分类直接对应，且三者的时间信号均早于物品域选择，与OFC在经济决策中的因果性地位相符\citep{ballesta2020values,ballesta2022orbitofrontal}。本研究的核心发现之一是这一价值编码框架在财富动态变化的背景下依然成立。

在财富编码层面，本研究在单神经元水平发现了OFC对当前财富状态的显著表征，与Juechems等人\citep{juechems2017ventromedial}在人类fMRI研究中的发现一致。Nguyen等人\citep{nguyen2025neural}对猕猴额叶多区域进行了系统记录，却未在OFC发现与财富相关的显著编码，这与本课题的结果存在明显矛盾，这一分歧的具体原因有待进一步分析。

除了价值计算，OFC的功能可能还包括对当前任务状态空间的抽象表征\citep{wilson2014orbitofrontal,wikenheiser2016cognitive,knudsen2022taking}。Schuck等人\citep{schuck2016human}在人类fMRI研究中发现，从OFC的信号中能够解码出需从部分观察中推断出的不可见任务状态；与之对应，Juechems等人\citep{juechems2017ventromedial}研究中的财富水平同样以隐式方式存在，需受试者从试次历史中主动推断——两者共同指向一个反复浮现的关键变量：状态的可见性（visibility）。若要检验OFC是否将财富作为抽象的任务状态变量编码，则需要在实验中至少部分地隐藏财富状态，使得动物必须从试次历史中主动推断当前财富水平。进一步地，在当前任务中财富水平随赚取和消费单调变化，缺乏足够的状态转移多样性。若能在实验设计中控制财富状态的转移结构——例如取消每次代币变化均为1枚的限制，允许在一次选择中赚取或消费多枚代币，使财富状态按非连续的路径跳变——则可以直接检验OFC是否真正维持了一张财富状态的转移地图，而非仅仅追踪财富的连续变化。

\subsection{ACC}\label{chap4_2_10:ACC}

本研究发现ACC对财富状态的显著编码及动态更新，这与此前多项研究的方向一致\citep{seo2009behavioral,yang2022primate,ferro2026accumulation}，但与Nguyen等人\citep{nguyen2025neural}的结果存在细微区别：他们发现猕猴dACC编码的是财富水平依赖的预期价值，而非财富水平本身，财富状态的表征被定位在vACC，这一差异的具体原因有待进一步分析。

在价值编码层面，已有研究尚未给出一致的图景。多数工作报告ACC编码选项价值\citep{kennerley2009neurons, prevost2010separate, hosokawa2013mechanisms, nguyen2025neural, ferro2026accumulation}，但亦有阴性结果\citep{cai2012neuronal}。本研究进一步揭示出更为复杂的编码模式：在两类选项中，ACC仅编码消费选项的价值（$V_S$），而不编码赚取选项的价值（$V_E$）。这表明ACC在价值计算中的功能可能比通常设想的更为复杂，其具体角色仍有待澄清。

在选择编码层面，虽然ACC的物品域选择信号强于动作域信号，但若与DLPFC和OFC横向比较，ACC的物品域选择编码并不占据优势——即使DLPFC被经典研究定性为动作域编码脑区，其物品域信号也与ACC大致相当。这说明无论是物品域还是动作域信号，可能都不是ACC编码的主要内容。

综合来看，ACC在“具体”维度上并无强表征：价值层面只见消费价值而无赚取价值，选择层面物品域与动作域信号均不突出；真正被显著编码的，是不对应任何具体选项或动作的财富状态。这或许并非巧合：财富不绑定于单一选项或动作，而是刻画决策整体情境的高阶变量，正是在此意义上它是“抽象”的。ACC对其的强编码，与既有研究将ACC视为整合上下文、编码抽象状态变量的观点一致\citep{heilbronner2016dorsal, monosov2020anterior}。这一抽象状态编码与OFC的认知地图，或许是同一功能的不同侧面。若能在实验中控制财富状态的转移结构（参见~\ref{chap4_2_9:OFC}~节），或可检验ACC是否同样维持着一张财富状态的转移地图，并进一步辨析ACC与OFC在抽象状态表征上的分工。

\subsection{DLPFC}\label{chap4_2_11:DLPFC}

本研究中DLPFC以动作域信号为主，但同时编码了物品域信号，且物品域信号的时间晚于动作域信号。如第~\ref{chap4_2_2:serial_parallel}节的讨论，这可能揭示了一个复杂的串行并行交织的决策系统。这里的结果与Cai等人\citep{cai2014contributions}的经典工作有一个关键区别需要指出：在他们的实验中，物品信息与位置信息是先后呈现的，其结果只能说明DLPFC具备编码物品域信号的能力，而无法直接回答时序问题。Lin等人\citep{lin2020evidence}在同时呈现两类信息的条件下发现，DLPFC中物品域信号不早于动作域信号；本研究中DLPFC物品域信号时间晚于动作域信号，与Lin等的结论方向一致。这共同指向一个与串行转换框架相悖的图景：物品域信号并非DLPFC决策计算的起点，而是动作域信号出现之后才形成的。综合来看，DLPFC仍然在视觉空间框架中形成决策。

另外，当前任务设计中，每个消费选项的代币成本固定为1，每个赚取选项的单次代币收益也固定为1，导致选项的总价值与总成本之间高度共线，无法独立操控两者。引入价格与数量均可变化的选项设计——例如允许消费选项在饮用水量和代币成本上均独立变化，使得高水量高成本与低水量低成本之间的选择成为可能——将制造出总价值与总成本方向冲突的情境。在此类冲突情境中，"选择当前最具诱惑力的选项"与"维持长期资产的理性管理"之间的冲突显著增加，将有助于检验DLPFC在价值决策中的认知控制功能。

\subsection{环路层面的神经机制}\label{chap4_2_12:circuit}

受实验中所采用的记录工具的限制，我们未能在同一个实验日中稳定记录足够数量的高质量单神经元。故本课题所有神经分析均建立在伪群体（pseudo-population）数据之上，我们无法精确地在单个试次水平观察脑区之间的协同机制，例如判断脑区间的信息流向。未来研究可以通过引入高通量记录技术（如NeuroPixel电极）来克服这一限制，直接在单试次水平上分析不同脑区之间的动态交互，从而更深入地理解财富表征与价值计算的神经环路机制。

在此基础上，电刺激和光遗传学可以提供脑区必要性与充分性的证据；但此类操控的特异性受到神经元混合编码的限制，实验设计需对操控时机、强度和靶区进行严格规划，以最大程度地将目标信号与混淆信号解耦。

\subsection{行为模型的潜在局限与改进方向}\label{chap4_2_13:model}

在当前任务中，行为模型依赖两项核心简化假设：（1）饮用水奖励的主观价值不随时间和累积给水量变化；（2）盯点成本与注视时长呈线性关系。尽管该模型能较好地拟合猕猴的选择行为，这些假设可能存在不可忽视的局限。

首先，奖励价值可能展现出时间依赖性与长期累积效应。本模型将单位饮用水的价值视为常数，但随着实验的推进，猕猴持续获得奖励，其体内水分状态与饮水动机必然发生变化。行为上，猕猴的无差异点并未随累积给水量线性漂移，而是在单次记录接近尾声（约90\%进度）时突然出现动机崩溃，表现为拒绝任何选择。这提示水的价值函数可能不是连续变化，而更接近带阈值的非线性形式：当体内水分积累低于某一阈值时，边际价值相对稳定；一旦超过阈值，边际价值急剧下降，导致饮水动机迅速消退。为刻画这一关系，可将水视作一种类似财富的累积资源，并为之构建含阈值的价值函数。具体而言，可以在实验中特定时间允许动物自由摄入饮用水以系统操纵其累积水分状态，进而估计该函数的具体形态。这一设计可将水的边际价值变化与时间进程解耦，为建立更精确的决策模型提供参数依据。

其次，奖励价值还可能存在短期时间尺度效应。除整个实验日尺度的长期累积外，近期饮水历史也可能对当前水的价值产生局部调节：若最近若干试次集中获得大量奖励，当前价值可能被暂时压低，反之则升高。这种短期效应使试次间价值波动无法仅由全局累积水量解释，但其时间窗口尚不明确——可能为数秒、数十秒，也可能跨越数十个试次。要独立于时间分离这一效应，需借助特定实验操纵：通过控制不同时段奖励大小出现的频率，人为营造短期缺水或水分充足的情境，并确保这些情境的起始与结束同绝对实验进程解耦。

最后，盯点成本可能并非线性增长。当前模型假定注视成本与注视时长成正比，但长时间维持注视而不眨眼颇为困难，注视成本函数斜率随注视时长递增。实验中观察到，猕猴有时在赚取选择中途主动中断、短暂休息后继续，呈现“赚钱---中断---继续赚钱”的模式。这一模式难以用线性成本说明，却能被非线性成本自然解释：当累计成本超出某一界限后，短暂休息后重新开始可能优于持续坚持。若盯点成本确为非线性，决策函数中须引入非线性努力折扣项，否则将对选项主观价值产生错误估计，进而混淆价值编码的神经分析。

% ─────────────────────────────────────────────
\section{展望}\label{chap4_3:future}
% ─────────────────────────────────────────────

本研究在神经元层面揭示了财富状态如何嵌入前额叶的决策框架——它不是一个独立于决策系统之外的背景变量，而是实时调制价值编码的内部信号。这一发现的潜在意义超出了经济决策本身。前额叶是认知功能的核心基础，承载着从工作记忆、注意力调控到情绪调节、风险评估在内的广泛计算。若财富状态的神经表征深度嵌入这一系统，则财富水平的变化所引发的，可能不只是价值编码增益的调整，而是对整个前额叶计算环境的系统性重塑：低财富状态可能通过与认知控制环路的交互削弱对即时诱惑的抑制，通过与情绪调节系统的交互放大负性情绪对风险评估的干扰，或通过占用工作记忆资源降低多步推理的精度。这些影响的具体通路各异，但共同指向同一个更一般性的问题：财富水平在多大程度上、通过何种神经机制，系统性地塑造了人类认知与决策的整体质量。

在应用层面，上述框架开启了若干具体的可能性。例如，若财富表征可以被操控，则原则上可以通过神经反馈（neurofeedback）或非侵入性脑刺激等手段在不改变真实财富水平的条件下调整个体对财富状态的神经表征，为缓解财富匮乏状态对认知系统的负面影响提供新的干预手段。从猕猴电生理到上述应用场景之间存在相当的距离，但在神经机制层面为这些问题提供可检验的假说，或许正是神经经济学研究的价值所在。